%-----------------------------------------------------------
% TEMPLATE: Mathematical Derivation Draft
% PURPOSE: For documenting step-by-step derivations
% AUTHOR: [Your Name]
% DATE: \today
%-----------------------------------------------------------
\documentclass[11pt, oneside]{article}

%-------------------------
% 1. ESSENTIAL PACKAGES
%-------------------------
\usepackage{amsmath, amssymb, amsthm} % Core math typesetting
\usepackage{geometry}                 % Flexible page margins
\usepackage{microtype}                % Better typography
\usepackage{hyperref}                 % For internal references
\usepackage[shortlabels]{enumitem}    % Customizable lists
\usepackage{xcolor}                   % For colored text/comments
\usepackage{todonotes}                % For draft notes in margins
\usepackage{titling}                  % For custom title spacing
\usepackage{mathtools}                % Extends amsmath (for \coloneqq, etc.)
\usepackage{mdframed}                 % For framed boxes (theorems, notes)
% Reset equation counter for each lemma
\usepackage{etoolbox}
\allowdisplaybreaks
\AtBeginEnvironment{lemma}{\setcounter{equation}{0}}
%-------------------------
% 2. PAGE & GEOMETRY SETUP
%-------------------------
\geometry{
    a4paper,
    left=25mm,
    right=25mm,
    top=25mm,
    bottom=30mm,
    % Showframe  % Uncomment to see page margins during setup
}
\setlength{\parindent}{0pt}          % No paragraph indentation
\setlength{\parskip}{0.8em}          % Space between paragraphs
\setlength{\droptitle}{-4em}         % Adjust title position

%-------------------------
% 3. CUSTOM COLORS & COMMANDS
%-------------------------
\definecolor{deepblue}{RGB}{0,0,128}
\definecolor{deepgreen}{RGB}{0,100,0}
\definecolor{noteorange}{RGB}{200,100,0}
\hypersetup{colorlinks=true, linkcolor=deepblue, urlcolor=deepblue}

% Quick mathematical symbols/operators
\newcommand{\R}{\mathbb{R}}          % Real numbers
\newcommand{\C}{\mathbb{C}}          % Complex numbers
\newcommand{\diff}{\mathop{}\!\mathrm{d}} % Proper differential 'd'
\newcommand{\pdiff}[2]{\frac{\partial #1}{\partial #2}} % Partial derivative
\newcommand{\eqnref}[1]{Eq.~(\ref{#1})} % Clean equation reference

%-------------------------
% 4. CUSTOM ENVIRONMENTS
%-------------------------
% For stating assumptions explicitly
\newenvironment{assumptions}
    {\begin{mdframed}[linecolor=deepgreen, linewidth=1.5pt,
         frametitle={Assumptions}, frametitlealignment=\centering]}
    {\end{mdframed}}
\newenvironment{lemma}
    {\begin{mdframed}[linecolor=deepgreen, linewidth=1.5pt,
         frametitle={Lemma}, frametitlealignment=\centering]}
    {\end{mdframed}}
% For key results or final formulas
\newenvironment{result}
    {\begin{mdframed}[linecolor=deepblue, linewidth=2pt, frametitle={Result}, frametitlealignment=\centering]}
    {\end{mdframed}}

% For important derivation steps
\newenvironment{keystep}
    {\vspace{0.5em}\noindent\textbf{Key Step:} \itshape}
    {\normalfont\vspace{0.5em}}

% For placeholder notes (remove in final draft)
\newcommand{\derivnote}[1]{\todo[inline, color=noteorange!30]{#1}}

%-------------------------
% 5. DOCUMENT START
%-------------------------
\begin{document}
\section{Cumulant Expression}
\begin{multline}
    \lambda^{pqr}_{stu}=\gamma^{pqr}_{stu}- \gamma^{p}_{s} \gamma^{qr}_{tu} + \gamma^{p}_{t} \gamma^{qr}_{su} + \gamma^{p}_{u} \gamma^{qr}_{ts} + \gamma^{q}_{s} \gamma^{pr}_{tu} - \gamma^{q}_{t} \gamma^{pr}_{su} + \gamma^{q}_{u} \gamma^{pr}_{st} + \gamma^{r}_{s} \gamma^{qp}_{tu} + \gamma^{r}_{t} \gamma^{pq}_{su} - \gamma^{r}_{u} \gamma^{pq}_{st}\\
+2 \gamma^{p}_{s} \gamma^{q}_{t} \gamma^{r}_{u} - 2 \gamma^{p}_{s} \gamma^{q}_{u} \gamma^{r}_{t} - 2 \gamma^{p}_{t} \gamma^{q}_{s} \gamma^{r}_{u} + 2 \gamma^{p}_{t} \gamma^{q}_{u} \gamma^{r}_{s} + 2 \gamma^{p}_{u} \gamma^{q}_{s} \gamma^{r}_{t} - 2 \gamma^{p}_{u} \gamma^{q}_{t} \gamma^{r}_{s}
\end{multline}
\section{Normal Order}
\begin{multline}
    a^{pqr}_{stu}=\gamma^{p}_{s} \tilde{a}^{qr}_{tu} - \gamma^{p}_{t} \tilde{a}^{qr}_{su} - \gamma^{p}_{u} \tilde{a}^{qr}_{ts} - \gamma^{q}_{s} \tilde{a}^{pr}_{tu} + \gamma^{q}_{t} \tilde{a}^{pr}_{su} \\
- \gamma^{q}_{u} \tilde{a}^{pr}_{st} - \gamma^{r}_{s} \tilde{a}^{qp}_{tu} - \gamma^{r}_{t} \tilde{a}^{pq}_{su} + \gamma^{r}_{u} \tilde{a}^{pq}_{st} + \gamma^{pq}_{st} \tilde{a}^{r}_{u} \\
- \gamma^{pq}_{su} \tilde{a}^{r}_{t} - \gamma^{pq}_{ut} \tilde{a}^{r}_{s} - \gamma^{pr}_{st} \tilde{a}^{q}_{u} + \gamma^{pr}_{su} \tilde{a}^{q}_{t} - \gamma^{pr}_{tu} \tilde{a}^{q}_{s} \\
- \gamma^{qr}_{su} \tilde{a}^{p}_{t} + \gamma^{qr}_{tu} \tilde{a}^{p}_{s} - \gamma^{rq}_{st} \tilde{a}^{p}_{u} + \gamma^{pqr}_{stu} + \tilde{a}^{pqr}_{stu}\\

\end{multline}
\begin{equation}
    a^{pq}_{st}=\gamma^{p}_{s} \tilde{a}^{q}_{t} - \gamma^{p}_{t} \tilde{a}^{q}_{s} - \gamma^{q}_{s} \tilde{a}^{p}_{t} + \gamma^{q}_{t} \tilde{a}^{p}_{s} + \gamma^{pq}_{st} + \tilde{a}^{pq}_{st}
\end{equation}
\begin{equation}
    a^{p}_{s}=\gamma^{p}_{s} + \tilde{a}^{p}_{s}
\end{equation}
\section{Change to Operator}
\begin{equation}
    \tilde{a}^{p}_{s}=- \gamma^{p}_{s} + a^{p}_{s}
\end{equation}
\begin{equation}
    \tilde{a}^{pq}_{st}=- \gamma^{p}_{s} \tilde{a}^{q}_{t} + \gamma^{p}_{t} \tilde{a}^{q}_{s} + \gamma^{q}_{s} \tilde{a}^{p}_{t} - \gamma^{q}_{t} \tilde{a}^{p}_{s} - \gamma^{pq}_{st} + a^{pq}_{st}
\end{equation}
\section{Final Result}
\begin{multline}
 a^{pqr}_{stu}=6 \gamma^{p}_{s} \gamma^{q}_{t} \gamma^{r}_{u} - 2 \gamma^{p}_{s} \gamma^{q}_{t} a^{r}_{u} - 6 \gamma^{p}_{s} \gamma^{q}_{u} \gamma^{r}_{t} + 2 \gamma^{p}_{s} \gamma^{q}_{u} a^{r}_{t} + 2 \gamma^{p}_{s} \gamma^{r}_{t} a^{q}_{u} \\
- 2 \gamma^{p}_{s} \gamma^{r}_{u} a^{q}_{t} - 2 \gamma^{p}_{s} \gamma^{qr}_{tu} + \gamma^{p}_{s} a^{qr}_{tu} - 6 \gamma^{p}_{t} \gamma^{q}_{s} \gamma^{r}_{u} + 2 \gamma^{p}_{t} \gamma^{q}_{s} a^{r}_{u} \\
+ 6 \gamma^{p}_{t} \gamma^{q}_{u} \gamma^{r}_{s} - 2 \gamma^{p}_{t} \gamma^{q}_{u} a^{r}_{s} - 2 \gamma^{p}_{t} \gamma^{r}_{s} a^{q}_{u} + 2 \gamma^{p}_{t} \gamma^{r}_{u} a^{q}_{s} + 2 \gamma^{p}_{t} \gamma^{qr}_{su} \\
- \gamma^{p}_{t} a^{qr}_{su} + 6 \gamma^{p}_{u} \gamma^{q}_{s} \gamma^{r}_{t} - 2 \gamma^{p}_{u} \gamma^{q}_{s} a^{r}_{t} - 6 \gamma^{p}_{u} \gamma^{q}_{t} \gamma^{r}_{s} + 2 \gamma^{p}_{u} \gamma^{q}_{t} a^{r}_{s} \\
+ 2 \gamma^{p}_{u} \gamma^{r}_{s} a^{q}_{t} - 2 \gamma^{p}_{u} \gamma^{r}_{t} a^{q}_{s} + \gamma^{p}_{u} \gamma^{qr}_{ts} + \gamma^{p}_{u} \gamma^{rq}_{st} - \gamma^{p}_{u} a^{qr}_{ts} \\
- 2 \gamma^{q}_{s} \gamma^{r}_{t} a^{p}_{u} + 2 \gamma^{q}_{s} \gamma^{r}_{u} a^{p}_{t} + 2 \gamma^{q}_{s} \gamma^{pr}_{tu} - \gamma^{q}_{s} a^{pr}_{tu} + 2 \gamma^{q}_{t} \gamma^{r}_{s} a^{p}_{u} \\
- 2 \gamma^{q}_{t} \gamma^{r}_{u} a^{p}_{s} - 2 \gamma^{q}_{t} \gamma^{pr}_{su} + \gamma^{q}_{t} a^{pr}_{su} - 2 \gamma^{q}_{u} \gamma^{r}_{s} a^{p}_{t} + 2 \gamma^{q}_{u} \gamma^{r}_{t} a^{p}_{s} \\
+ 2 \gamma^{q}_{u} \gamma^{pr}_{st} - \gamma^{q}_{u} a^{pr}_{st} + \gamma^{r}_{s} \gamma^{pq}_{ut} + \gamma^{r}_{s} \gamma^{qp}_{tu} - \gamma^{r}_{s} a^{qp}_{tu} \\
+ 2 \gamma^{r}_{t} \gamma^{pq}_{su} - \gamma^{r}_{t} a^{pq}_{su} - 2 \gamma^{r}_{u} \gamma^{pq}_{st} + \gamma^{r}_{u} a^{pq}_{st} + \gamma^{pq}_{st} a^{r}_{u} \\
- \gamma^{pq}_{su} a^{r}_{t} - \gamma^{pq}_{ut} a^{r}_{s} - \gamma^{pr}_{st} a^{q}_{u} + \gamma^{pr}_{su} a^{q}_{t} - \gamma^{pr}_{tu} a^{q}_{s} \\
- \gamma^{qr}_{su} a^{p}_{t} + \gamma^{qr}_{tu} a^{p}_{s} - \gamma^{rq}_{st} a^{p}_{u} + \gamma^{pqr}_{stu} + \tilde{a}^{pqr}_{stu}\\

\end{multline}
\begin{multline}
    a^{pqr}_{stu}=\tilde{a}^{pqr}_{stu}+\gamma^{p}_{s} a^{qr}_{tu} - \gamma^{p}_{t} a^{qr}_{su} - \gamma^{p}_{u} a^{qr}_{ts} - \gamma^{q}_{s} a^{pr}_{tu} + \gamma^{q}_{t} a^{pr}_{su} - \gamma^{q}_{u} a^{pr}_{st} - \gamma^{r}_{s} a^{qp}_{tu} - \gamma^{r}_{t} a^{pq}_{su} + \gamma^{r}_{u} a^{pq}_{st}\\
+\gamma^{pq}_{st} a^{r}_{u} - \gamma^{pq}_{su} a^{r}_{t} - \gamma^{pq}_{ut} a^{r}_{s} - \gamma^{pr}_{st} a^{q}_{u} + \gamma^{pr}_{su} a^{q}_{t} - \gamma^{pr}_{tu} a^{q}_{s} - \gamma^{qr}_{su} a^{p}_{t} + \gamma^{qr}_{tu} a^{p}_{s} - \gamma^{rq}_{st} a^{p}_{u}\\
- 2 \gamma^{p}_{s} \gamma^{q}_{t} a^{r}_{u} - 2 \gamma^{p}_{s} \gamma^{r}_{u} a^{q}_{t} - 2 \gamma^{p}_{t} \gamma^{q}_{u} a^{r}_{s} - 2 \gamma^{p}_{t} \gamma^{r}_{s} a^{q}_{u} - 2 \gamma^{p}_{u} \gamma^{q}_{s} a^{r}_{t} - 2 \gamma^{p}_{u} \gamma^{r}_{t} a^{q}_{s} - 2 \gamma^{q}_{s} \gamma^{r}_{t} a^{p}_{u} - 2 \gamma^{q}_{t} \gamma^{r}_{u} a^{p}_{s} - 2 \gamma^{q}_{u} \gamma^{r}_{s} a^{p}_{t}\\
2 \gamma^{p}_{s} \gamma^{q}_{u} a^{r}_{t} + 2 \gamma^{p}_{s} \gamma^{r}_{t} a^{q}_{u} + 2 \gamma^{p}_{t} \gamma^{q}_{s} a^{r}_{u} + 2 \gamma^{p}_{t} \gamma^{r}_{u} a^{q}_{s} + 2 \gamma^{p}_{u} \gamma^{q}_{t} a^{r}_{s} + 2 \gamma^{p}_{u} \gamma^{r}_{s} a^{q}_{t} + 2 \gamma^{q}_{s} \gamma^{r}_{u} a^{p}_{t} + 2 \gamma^{q}_{t} \gamma^{r}_{s} a^{p}_{u} + 2 \gamma^{q}_{u} \gamma^{r}_{t} a^{p}_{s}\\
- \gamma^{p}_{s} \gamma^{qr}_{tu} + \gamma^{p}_{t} \gamma^{qr}_{su} + \gamma^{p}_{u} \gamma^{rq}_{st} + \gamma^{q}_{s} \gamma^{pr}_{tu} - \gamma^{q}_{t} \gamma^{pr}_{su} + \gamma^{q}_{u} \gamma^{pr}_{st} + \gamma^{r}_{s} \gamma^{pq}_{ut} + \gamma^{r}_{t} \gamma^{pq}_{su} - \gamma^{r}_{u} \gamma^{pq}_{st}\\
+4 \gamma^{p}_{s} \gamma^{q}_{t} \gamma^{r}_{u} - 4 \gamma^{p}_{s} \gamma^{q}_{u} \gamma^{r}_{t} - 4 \gamma^{p}_{t} \gamma^{q}_{s} \gamma^{r}_{u} + 4 \gamma^{p}_{t} \gamma^{q}_{u} \gamma^{r}_{s} + 4 \gamma^{p}_{u} \gamma^{q}_{s} \gamma^{r}_{t} - 4 \gamma^{p}_{u} \gamma^{q}_{t} \gamma^{r}_{s}+\lambda^{pqr}_{stu}
\end{multline}
\section{Spin Free Result}
\begin{multline}
    E^{pqr}_{stu}\approx 
D^{p}_{s} E^{qr}_{tu} - \frac{D^{p}_{t} E^{qr}_{su}}{2} - \frac{D^{p}_{u} E^{qr}_{ts}}{2} - \frac{D^{q}_{s} E^{pr}_{tu}}{2} + D^{q}_{t} E^{pr}_{su} - \frac{D^{q}_{u} E^{pr}_{st}}{2} - \frac{D^{r}_{s} E^{qp}_{tu}}{2} - \frac{D^{r}_{t} E^{pq}_{su}}{2} + D^{r}_{u} E^{pq}_{st}\\
+D^{pq}_{st} E^{r}_{u} - \frac{D^{pq}_{su} E^{r}_{t}}{2} - \frac{D^{pq}_{ut} E^{r}_{s}}{2} - \frac{D^{pr}_{st} E^{q}_{u}}{2} + D^{pr}_{su} E^{q}_{t} - \frac{D^{pr}_{tu} E^{q}_{s}}{2} - \frac{D^{qr}_{su} E^{p}_{t}}{2} + D^{qr}_{tu} E^{p}_{s} - \frac{D^{rq}_{st} E^{p}_{u}}{2}\\
- 2 D^{p}_{s} D^{q}_{t} E^{r}_{u} - 2 D^{p}_{s} D^{r}_{u} E^{q}_{t} - \frac{D^{p}_{t} D^{q}_{u} E^{r}_{s}}{2} - \frac{D^{p}_{t} D^{r}_{s} E^{q}_{u}}{2} - \frac{D^{p}_{u} D^{q}_{s} E^{r}_{t}}{2} - \frac{D^{p}_{u} D^{r}_{t} E^{q}_{s}}{2} - \frac{D^{q}_{s} D^{r}_{t} E^{p}_{u}}{2} - 2 D^{q}_{t} D^{r}_{u} E^{p}_{s} - \frac{D^{q}_{u} D^{r}_{s} E^{p}_{t}}{2}\\
+D^{p}_{s} D^{q}_{u} E^{r}_{t} + D^{p}_{s} D^{r}_{t} E^{q}_{u} + D^{p}_{t} D^{q}_{s} E^{r}_{u} + D^{p}_{t} D^{r}_{u} E^{q}_{s} + D^{p}_{u} D^{q}_{t} E^{r}_{s} + D^{p}_{u} D^{r}_{s} E^{q}_{t} + D^{q}_{s} D^{r}_{u} E^{p}_{t} + D^{q}_{t} D^{r}_{s} E^{p}_{u} + D^{q}_{u} D^{r}_{t} E^{p}_{s}\\
- D^{p}_{s} D^{qr}_{tu} + \frac{D^{p}_{t} D^{qr}_{su}}{2} + \frac{D^{p}_{u} D^{rq}_{st}}{2} + \frac{D^{q}_{s} D^{pr}_{tu}}{2} - D^{q}_{t} D^{pr}_{su} + \frac{D^{q}_{u} D^{pr}_{st}}{2} + \frac{D^{r}_{s} D^{pq}_{ut}}{2} + \frac{D^{r}_{t} D^{pq}_{su}}{2} - D^{r}_{u} D^{pq}_{st}\\
+4 D^{p}_{s} D^{q}_{t} D^{r}_{u} - 2 D^{p}_{s} D^{q}_{u} D^{r}_{t} - 2 D^{p}_{t} D^{q}_{s} D^{r}_{u} + D^{p}_{t} D^{q}_{u} D^{r}_{s} + D^{p}_{u} D^{q}_{s} D^{r}_{t} - 2 D^{p}_{u} D^{q}_{t} D^{r}_{s}\\

\end{multline}
\end{document}